\paragraph{Lasso Regression}

\subparagraph{Hạng tử điều chỉnh}
Khác với Ridge Regression sử dụng chuẩn $L_2$, Lasso Regression sử dụng chuẩn
$L_1$ cho vector trọng số. Hạng tử điều chỉnh có dạng:
\begin{equation}
  E_W(\boldsymbol{w}) = \|\boldsymbol{w}\|_1
\end{equation}
Hạng tử này là tổng giá trị tuyệt đối của các trọng số. Việc sử dụng chuẩn
$L_1$ cũng nhằm ``phạt'' các trọng số lớn, tuy nhiên theo cách khác với Ridge.

\subparagraph{Hàm mục tiêu}
Hàm mục tiêu của Lasso Regression được viết như sau:
\begin{equation}
  E(\boldsymbol{w}) = \frac{1}{2} \|\boldsymbol{X}\boldsymbol{w} -
  \boldsymbol{t}\|^2 + \lambda \|\boldsymbol{w}\|_1.
\end{equation}
\subparagraph{Nghiệm giải}
Không giống như Ridge Regression, Lasso Regression
\textbf{không có nghiệm đóng dạng tường minh}. Nguyên nhân là do hàm $|w|$
không khả vi tại $w = 0$. Do đó, nghiệm của Lasso thường được tìm bằng các
phương pháp tối ưu số, chẳng hạn như:
\begin{itemize}
  \item Gradient Descent (với subgradient)
  \item Coordinate Descent
  \item Least Angle Regression (LARS)
\end{itemize}

\subparagraph{Đặc điểm}
Một tính chất quan trọng của Lasso Regression là khả năng tạo ra nghiệm thưa
(``sparse solution''), tức là nhiều trọng số có giá trị đúng bằng 0.
\begin{itemize}
  \item Ridge Regression: làm nhỏ các trọng số nhưng hiếm khi bằng 0
  \item Lasso Regression: có thể đưa một số trọng số về đúng 0
\end{itemize}
Nhờ đó, Lasso có thể được sử dụng như một phương pháp
\textbf{lựa chọn đặc trưng} (feature selection).

\subparagraph{Diễn giải trực quan}
Có thể hiểu Lasso Regression là bài toán tối ưu với ràng buộc:
\begin{equation}
  \min_{\boldsymbol{w}} \|\boldsymbol{X}\boldsymbol{w} - \boldsymbol{t}\|^2
  \quad \text{với} \quad\|\boldsymbol{w}\|_1 \leq C.
\end{equation}

Miền ràng buộc $\|\boldsymbol{w}\|_1 \leq C$ có dạng hình thoi (diamond). Khi
kết hợp với các đường đồng mức của hàm lỗi, nghiệm tối ưu thường rơi vào các
đỉnh của hình thoi, dẫn đến nhiều hệ số bằng 0.
